\documentclass[12pt]{article}

\usepackage{sbc-template}
\usepackage{graphicx,url}
\usepackage[utf8]{inputenc}
\usepackage[T1]{fontenc}
\usepackage[brazil]{babel}

\sloppy

\title{Ideia do padrão decorator em um sistema de hamburgueria}

\author{Raul Beltrame\inst{1}, Matheus Dias\inst{1}}

\address{Curso de Programação Orientada a Objetos\\
Sistemas de Informações - UFVJM\\
\email{\{raul.beltrame,matheus.dias\}@ufvjm.edu.br}
}

\begin{document}

\maketitle

\begin{resumo}
Neste artigo analisamos qual padrão de projeto era o mais simples e coerente para aplicarmos
em um sistema de hamburgueria desenvolvido em Java na disciplina programação
orientada a objetos. Neste estudo comparamos padrões candidatos com a estrutura atual
considerando esforço de implementação, impacto nas classes
que já temos e também se tinha uma certa clareza didática. Nos estudos pensamos que
o padrão Decorator era a alternativa mais adequada para o projeto, pois o sistema
já possui produtos, adicionais e itens de pedido que preço final depende de opcionais.
\end{resumo}

\section{Introdução}

Os Padrões de projeto são soluções reutilizáveis para problemas recorrentes de
desenvolvimento de software orientado a objetos. Eles não representam exatamente código
pronto, mas formas de colaboração entre classes que ajudam a organizar
responsabilidades, reduzir acoplamento e melhorar a manutenção do sistema
\cite{gamma1994}.

Essa preocupação também se relaciona com critérios de qualidade de
software que venho estudando. Kleppmann destaca confiabilidade, escalabilidade e manutenibilidade
como objetivos importantes na construção de sistemas, além de associar a
manutenibilidade a qualidades como simplicidade, operabilidade e
facilidade de evolução \cite{kleppmann2017}. Mesmo em um projeto de menor escala,
esses conceitos ajudam a justificar o uso de padrões de projeto,
pois uma solução bem organizada tende a ser mais fácil de entender, reusar, alterar e 
até para nossa apresentação.

Falando do projeto, o projeto analisado é um sistema de hamburgueria desenvolvido em Java. Ele
possui classes para clientes, funcionários, produtos, pedidos, vendas,
estoque, entregas e persistência em JSON. Dentro desse escopo, uma regra
importante é a montagem de um pedido a partir de um produto base e de
adicionais escolhidos pelo cliente. Essa característica aproxima o sistema do
padrão Decorator, cujo objetivo é acrescentar responsabilidades ou
comportamentos a um objeto de forma flexível, sem criar uma grande hierarquia
de subclasses.

Entre os padrões listados como opções, também consideramos o Singleton. Porém, considerando o
contexto da hamburgueria, achamos que o Decorator apresenta a melhor relação entre
facilidade de implementação, utilidade no domínio e clareza para explicação
acadêmica. Além disso, ele nos pareceu mais desafiador do que o Singleton.

\section{Objetivo}

O objetivo deste trabalho é identificar o padrão de projeto mais
adequado para ser implementado no sistema de hamburgueria, alinhando a
escolha com base na estrutura atual do código.
Especificamente, buscaamos avaliar o uso do padrão Decorator na
representação de produtos que já eram personalizados/decorados com adicionais, 
demonstrando como ele pode melhorar a organização da regra de cálculo do subtotal 
de um item de pedido.

\section{Material e Métodos}

O material que trabalhamos em cima foi o código do sistema de hamburgueria, estruturado
como um projeto em Java. A análise concentrou-se nas nossas classes
\texttt{Produto}, \texttt{DescricaoProduto}, \texttt{Adicional},
\texttt{ItemPedido}, \texttt{Pedido}, \texttt{CadastroPedidos},
\texttt{GeradorId} e \texttt{PersistenciaJSON}.

O método adotado foi uma análise  qualitativa do código. Para cada padrão
candidato, foram considerados quatro ideias:

\begin{itemize}
    \item esforço para implementar
    \item quantidade de classes a mudar
    \item coerência com as regras de negócio
    \item facilidade de explicação no artigo e na apresentação
\end{itemize}

\section{Resultados e Discussão}

Nosso estudo mostrou que o padrão Singleton seria o mais curto de implementar,
pois a classe \texttt{GeradorId} já centraliza a geração de identificadores por
meio de atributos e métodos estáticos. Bastaria transformar essa classe em uma
instância única controlada. Porém, essa solução teria impacto funcional
limitado, pois apenas formalizaria um comportamento que já está centralizado no
projeto.

O padrão Decorator, por outro lado, encaixa diretamente na regra de negócio
dos pedidos. Atualmente, a classe \texttt{ItemPedido} possui um \texttt{Produto}
e uma lista de objetos \texttt{Adicional}. O método de cálculo do subtotal soma
o preço base do produto com os valores dos adicionais e multiplica o resultado
pela quantidade. Essa estrutura já representa, de forma simplificada, a ideia
central do Decorator que é a partir de um objeto base e acrescentar funcionalidades ou
valores sem modificar diretamente a classe original.

Uma implementação formal poderia criar uma interface, por exemplo
\texttt{ItemCardapio}, com métodos para obter descrição e valor. O produto base
implementaria essa interface, enquanto cada adicional seria representado por um
decorador que também implementa a interface e mantém uma referência para outro
\texttt{ItemCardapio}. Assim, um hambúrguer poderia ser decorado com bacon,
queijo, molho extra ou outros adicionais de forma encadeada.

\begin{table}[ht]
\centering
\caption{Comparação dos padrões mais viáveis para o projeto}
\label{tab:comparacao}
\begin{tabular}{p{3.2cm}p{7.6cm}p{2.7cm}}
\hline
\textbf{Padrão} & \textbf{Aplicação no projeto} & \textbf{Avaliação} \\
\hline
Singleton & Geração única de identificadores na classe \texttt{GeradorId}. & Muito simples, mas pouco expressivo. \\
Decorator & Personalização de produtos com adicionais em \texttt{ItemPedido}. & Simples e muito aderente ao domínio. \\
Factory Method/Builder & Criação de pedidos e produtos complexos. & Possível, mas menos direto que Decorator. \\
\hline
\end{tabular}
\end{table}

Com base nessa comparação, o Decorator foi considerado por nós o padrão mais adequado
para o trabalho. Ele não exige mudanças profundas na arquitetura, aproveita
classes já existentes e melhora a separação de responsabilidades. A classe
\texttt{Produto} pode continuar representando o item cadastrado no cardápio,
enquanto os adicionais passam a ser tratados como elementos que decoram o item
base durante o cálculo do preço e da descrição final.

Outro ponto positivo é a clareza didática. Em uma hamburgueria, é natural
explicar que um lanche base recebe complementos opcionais. Isso facilita a
apresentação do padrão e mostra uma aplicação prática, diferente de uma
implementação meramente artificial.

Os demais padrões da lista também poderiam ser aplicados em
partes específicas do sistema. Porém, para este projeto, eles exigiriam a
criação de regras novas ou uma remodelagem maior do código, o que achamos
menos indicados para uma implementação simples e objetiva.

\section{Conclusão}

Concluimos que o padrão Decorator é a melhor escolha para o contexto do sistema
de hamburgueria. Embora o Singleton seja tecnicamente o padrão com menor
alteração, o Decorator apresenta melhor aderência ao escopo, pois
representa diretamente a personalização de produtos com adicionais, como já abstraímos.

A implementação proposta permite calcular o valor de um item de pedido de
forma mais flexível e organizada, sem criar subclasses para cada combinação de
hambúrguer e adicional. Dessa forma, o padrão contribui tanto para a qualidade
do código quanto para a clareza da explicação acadêmica do projeto.

\renewcommand{\refname}{Referências Bibliográficas}
\begin{thebibliography}{99}

\bibitem{gamma1994}
Gamma, E., Helm, R., Johnson, R. and Vlissides, J. (1994).
\newblock \textit{Design Patterns: Elements of Reusable Object-Oriented Software}.
\newblock Addison-Wesley.

\bibitem{freeman2004}
Freeman, E., Freeman, E., Sierra, K. and Bates, B. (2004).
\newblock \textit{Head First Design Patterns}.
\newblock O'Reilly Media.

\bibitem{kleppmann2017}
Kleppmann, M. (2017).
\newblock \textit{Designing Data-Intensive Applications: The Big Ideas Behind Reliable, Scalable, and Maintainable Systems}.
\newblock O'Reilly Media.

\bibitem{sbc2017}
Sociedade Brasileira de Computação (2017).
\newblock \textit{SBC Conferences Template -- updated sbc-template.sty v2017}.
\newblock Disponível em:
\url{https://www.overleaf.com/latex/templates/sbc-conferences-template-updated-sbc-template-dot-sty-v2017/pyhttxftxjqn}.
\newblock Acesso em: 05 maio 2026.

\end{thebibliography}

\end{document}
